
%(BEGIN_QUESTION)
% Copyright 2010, Tony R. Kuphaldt, released under the Creative Commons Attribution License (v 1.0)
% This means you may do almost anything with this work of mine, so long as you give me proper credit

Les National Transportation Safety Boards ulykkesnotat for rørledning (DCA-97-FP-002) om dieselrørbruddet i Murfreesboro, Tennessee i 1996, og svar på spørsmålene under.

\vskip 10pt

Identifiser den/de fysiske årsakene til rørbruddet: hva gjorde konkret at rørtrykket steg til punktet der ledningen rupturerte?

\vskip 10pt

En medvirkende faktor i ulykken var et feilaktig skjema på SCADA-skjermen for pumpestasjonen i Murfreesboro. Beskriv hva som var feil i skjemaet, og hvordan denne feilen villedet operatøren som overvåket det.

\vskip 10pt

Hvorfor nektet den elektriske blokkeringsventilen i Murfreesboro å åpne da operatøren ga åpne-kommando kl. 09:35:02? Basert på kunnskap om ulike ventiltyper: forklar hvorfor en slik ventil ble påvirket på denne måten.

\vskip 10pt

En tekniker som ble tilkalt for å kontrollere den elektriske blokkeringsventilen og tilhørende styreutstyr i Murfreesboro rapporterte at det «ikke var noen problemer». Hva mener du om denne vurderingen, gitt det du vet om ventilen og hvorfor den ikke åpnet? Etter ditt syn, ble ventilen kontrollert grundig nok?

\vskip 20pt \vbox{\hrule \hbox{\strut \vrule{} {\bf Forslag til sokratisk diskusjon} \vrule} \hrule}

\begin{itemize}
\item{} Side 3 i rapporten nevner en «depressurization wave» som brukte 74 sekunder fra bruddstedet til Coalmont. Hva tror du en slik «depressurization wave» er, og hvordan henger dette sammen med Pascals prinsipp?
\item{} Identifiser måter en instrumenttekniker kunne bidratt til å forhindre dette rørbruddet.
\end{itemize}

\underbar{file i00785}
%(END_QUESTION)





%(BEGIN_ANSWER)


%(END_ANSWER)





%(BEGIN_NOTES)

5. november 1996 forberedte Colonial Pipeline Company «pigging» av en dieselrørledning. Pumpene ved Atlanta-anlegget (oppstrøms) ble stoppet, og deretter ble en blokkeringsventil i Murfreesboro (nedstrøms) stengt. Piggingen skulle utføres på strekningen mellom Murfreesboro og en nedstrøms stasjon i Nashville. (Side 1)

\vskip 10pt

Kort tid etter at blokkeringsventilen i Murfreesboro var stengt, ble planen endret og pumpene oppstrøms den stengte ventilen startet igjen (i strid med Colonials prosedyrer). Dette førte til overtrykk i ledningen mellom pumpestasjonen i Coalmont (første stasjon oppstrøms Murfreesboro) og Murfreesboro. Overtrykket ble registrert i Coalmont, men feil innstilte trip-punkter i styringen gjorde at rørledningstrykket fikk overstige sikre grenser. (Side 2)

Trykktransmitteren i Murfreesboro registrerte ikke denne tilstanden fordi den var installert nedstrøms den stengte linjeblokkeringsventilen. Uheldigvis viste SCADA-grafikken transmitteren som om den var plassert oppstrøms ventilen, slik den var på de fleste andre stasjoner på denne linjen. (Side 3) Mer presist var forskjellen mellom Murfreesboro og øvrige stasjoner egentlig plasseringen av den elektriske linjeblokkeringsventilen, ikke transmitteren. På de fleste stasjoner går strømningen gjennom manuell blokkeringsventil, forbi transmitter, gjennom tilbakeslagsventil og deretter gjennom elektrisk blokkeringsventil. I Murfreesboro møtte strømningen først den elektriske blokkeringsventilen, så transmitteren, deretter tilbakeslagsventil og til slutt manuell blokkeringsventil. (Side 4)

\vskip 10pt

Seksten minutter etter pumpestart inn i den stengte linjen forsøkte operatøren første gang å åpne den stengte blokkeringsventilen i Murfreesboro (uten å stoppe pumpene, igjen i strid med prosedyrene). Ventilen åpnet ikke fordi differansetrykket over ventilen var for høyt (1700 PSI, mens aktuatoren bare var dimensjonert for 1440 PSI). Operatøren ringte en tekniker i Murfreesboro for kontroll av ventilen, men teknikeren rapporterte at den var i orden. (Side 3)

\vskip 10pt

Fire minutter senere stoppet regulatoren noen pumper og omdirigerte strømning ved en annen stasjon for å senke rørtrykket slik at Murfreesboro-ventilen kunne åpne. Den åpnet til slutt på fjerde forsøk, men først omtrent tre minutter etter at rørledningen allerede hadde rupturert på grunn av overtrykk. Et plutselig trykkfall ($\Delta P$ = 416 PSI) ble registrert i Coalmont etter bruddet, men SCADA-systemet genererte ingen alarmer. Estimert rørtrykk ved den rupturerte sveisen var 1820 PSI. «Maximum Operating Pressure» (MOP) for ledningen var 1318 PSI.

Operatøren fortsatte å pumpe diesel inn i den rupturerte ledningen i omtrent én time før problemet ble oppdaget, fordi forventet trykkøkning uteble.

\vfil \eject

\noindent
{\bf Forberedelsesquiz:}

En medvirkende faktor i ulykken var feilaktig skjema i SCADA-visningen for pumpestasjonen i Murfreesboro. Beskriv hva som var feil i skjemaet, og hvordan dette villedet operatøren.

\begin{itemize}
\item{} Trykket var feil vist i PSI i stedet for enheten «bar»
\vskip 5pt
\item{} Trykktransmitteren var feilaktig vist oppstrøms blokkeringsventilen
\vskip 5pt
\item{} Komposisjonsvisningen viste feilaktig bensin i stedet for diesel i linjen
\vskip 5pt
\item{} Blokkeringsventilen var feilaktig vist som luft-til-åpne i stedet for luft-til-lukke
\vskip 5pt
\item{} Pumpen var feilaktig vist med motsatt pumperetning
\vskip 5pt
\item{} Blokkeringsventilen var feilaktig vist som luft-til-lukke i stedet for luft-til-åpne
\end{itemize}


%INDEX% Leseoppgave: NTSB-rapport om dieselrørbruddet i Murfreesboro, Tennessee i 1996

%(END_NOTES)


